% Source: https://artofproblemsolving.com/community/c3h2745293_divisibility_rule_modular_arithmetic
% Stable identity: aops-topic-2745293-q01
\textbf{(a)} Each test has the form $p\mid a\iff p\mid b+kD$, where $10k\equiv1\pmod p$:
\[10\cdot(-2)=-20\equiv1\pmod7,\quad10\cdot4=40\equiv1\pmod{13},\quad10\cdot(-5)=-50\equiv1\pmod{17},\]
\[10\cdot2=20\equiv1\pmod{19},\quad10\cdot7=70\equiv1\pmod{23}.\]
Then
\[k\,a=10k\,b+kD\equiv b+kD\pmod p.\]
Since $p$ is prime and $p\nmid k$, $p\mid a$ if and only if $p\mid ka$, i.e. if and only if $p\mid b+kD$.

\textbf{(b)} Since $10\equiv-1\pmod{11}$,
\[a=\sum_{j=1}^{n}D_j\,10^{\,n-j}\equiv\sum_{j=1}^{n}D_j(-1)^{n-j}=(-1)^{n+1}\left(D_1-D_2+D_3-\cdots+(-1)^{n+1}D_n\right)\pmod{11}.\]
So $11\mid a$ exactly when $11$ divides the alternating sum. $\blacksquare$
